% Options for packages loaded elsewhere
\PassOptionsToPackage{unicode}{hyperref}
\PassOptionsToPackage{hyphens}{url}
\documentclass[
]{article}
\usepackage{xcolor}
\usepackage[margin=1in]{geometry}
\usepackage{amsmath,amssymb}
\setcounter{secnumdepth}{-\maxdimen} % remove section numbering
\usepackage{iftex}
\ifPDFTeX
  \usepackage[T1]{fontenc}
  \usepackage[utf8]{inputenc}
  \usepackage{textcomp} % provide euro and other symbols
\else % if luatex or xetex
  \usepackage{unicode-math} % this also loads fontspec
  \defaultfontfeatures{Scale=MatchLowercase}
  \defaultfontfeatures[\rmfamily]{Ligatures=TeX,Scale=1}
\fi
\usepackage{lmodern}
\ifPDFTeX\else
  % xetex/luatex font selection
\fi
% Use upquote if available, for straight quotes in verbatim environments
\IfFileExists{upquote.sty}{\usepackage{upquote}}{}
\IfFileExists{microtype.sty}{% use microtype if available
  \usepackage[]{microtype}
  \UseMicrotypeSet[protrusion]{basicmath} % disable protrusion for tt fonts
}{}
\makeatletter
\@ifundefined{KOMAClassName}{% if non-KOMA class
  \IfFileExists{parskip.sty}{%
    \usepackage{parskip}
  }{% else
    \setlength{\parindent}{0pt}
    \setlength{\parskip}{6pt plus 2pt minus 1pt}}
}{% if KOMA class
  \KOMAoptions{parskip=half}}
\makeatother
\usepackage{graphicx}
\makeatletter
\newsavebox\pandoc@box
\newcommand*\pandocbounded[1]{% scales image to fit in text height/width
  \sbox\pandoc@box{#1}%
  \Gscale@div\@tempa{\textheight}{\dimexpr\ht\pandoc@box+\dp\pandoc@box\relax}%
  \Gscale@div\@tempb{\linewidth}{\wd\pandoc@box}%
  \ifdim\@tempb\p@<\@tempa\p@\let\@tempa\@tempb\fi% select the smaller of both
  \ifdim\@tempa\p@<\p@\scalebox{\@tempa}{\usebox\pandoc@box}%
  \else\usebox{\pandoc@box}%
  \fi%
}
% Set default figure placement to htbp
\def\fps@figure{htbp}
\makeatother
\setlength{\emergencystretch}{3em} % prevent overfull lines
\providecommand{\tightlist}{%
  \setlength{\itemsep}{0pt}\setlength{\parskip}{0pt}}
\usepackage{bookmark}
\IfFileExists{xurl.sty}{\usepackage{xurl}}{} % add URL line breaks if available
\urlstyle{same}
\hypersetup{
  hidelinks,
  pdfcreator={LaTeX via pandoc}}

\author{}
\date{\vspace{-2.5em}}

\begin{document}

\section{============================================================}\label{section}

\section{Analyse Multivariée du Bonheur
Mondial}\label{analyse-multivariuxe9e-du-bonheur-mondial}

\section{ACP, K-Means, CAH}\label{acp-k-means-cah}

\section{============================================================}\label{section-1}

tinytex::tlmgr\_install(``bookmark'') \# 1. Chargement des packages
--------------------------------- library(tidyverse) library(FactoMineR)
library(factoextra) library(cluster) library(corrplot)
library(gridExtra) library(dendextend) library(knitr)
library(kableExtra)

\section{2. Import des données
--------------------------------------}\label{import-des-donnuxe9es}

df \textless- read.csv(``C:/Users/ichra/Downloads/happiness\_data.csv'',
stringsAsFactors = FALSE) cat(``Dimensions du dataset :'', nrow(df),
``lignes x'', ncol(df), ``colonnes\n'')

\section{3. Variables numériques
------------------------------------}\label{variables-numuxe9riques}

num\_vars \textless- c(``Happiness\_Score'', ``GDP\_per\_Capita'',
``Social\_Support'', ``Healthy\_Life\_Expectancy'', ``Freedom'',
``Generosity'', ``Corruption\_Perception'', ``Unemployment\_Rate'',
``Education\_Index'')

\section{4. Description des variables (tableau)
--------------------}\label{description-des-variables-tableau}

desc\_vars \textless- data.frame( Variable = c(``Happiness\_Score'',
``GDP\_per\_Capita'', ``Social\_Support'',
``Healthy\_Life\_Expectancy'', ``Freedom'', ``Generosity'',
``Corruption\_Perception'', ``Unemployment\_Rate'',
``Education\_Index'', ``Region''), Type = c(rep(``Numérique'', 9),
``Catégorielle''), Description = c( ``Score de bonheur perçu (échelle
Cantril)'', ``Log du PIB par habitant (USD)'', ``Support social perçu
(0-1)'', ``Espérance de vie en bonne santé (années)'', ``Liberté de
choix de vie (0-1)'', ``Générosité perçue (dons)'', ``Perception faible
corruption (0-1)'', ``Taux de chômage (\%)'', ``Indice de niveau
d'éducation (0-1)'', ``Région géographique (7 modalités)'' ), Plage =
c(``2.0 -- 8.0'', ``6.5 -- 12.5'', ``0.30 -- 0.99'', ``45 -- 80'',
``0.20 -- 0.99'', ``-0.30 -- 0.60'', ``0.02 -- 0.55'', ``1.6 -- 18.8'',
``0.20 -- 0.99'', ``--'') ) kable(desc\_vars, caption = ``Description
des variables du dataset'', booktabs = TRUE, linesep = ``\,``)
\%\textgreater\% kable\_styling(latex\_options = c(''striped'',
``hold\_position''), font\_size = 10) \%\textgreater\% column\_spec(1,
bold = TRUE, width = ``4cm'') \%\textgreater\% column\_spec(3, width =
``6cm'') cat(``\n\n'')

\section{5. Distribution régionale
---------------------------------}\label{distribution-ruxe9gionale}

region\_tab \textless- as.data.frame(table(df\(Region))
colnames(region_tab) <- c("Région", "Nombre de pays")
region_tab\)Pourcentage \textless- round(100 *
region\_tab\$\texttt{Nombre\ de\ pays} / nrow(df), 1) kable(region\_tab,
caption = ``Distribution des pays par région'', booktabs = TRUE)
\%\textgreater\% kable\_styling(latex\_options = c(``striped'',
``hold\_position'')) cat(``\n\n'')

\section{6. Qualité des données
------------------------------------}\label{qualituxe9-des-donnuxe9es}

cat(``=== Valeurs manquantes ===\n'') print(colSums(is.na(df{[},
num\_vars{]}))) cat(``\nNombre de doublons :'', sum(duplicated(df)),
``\n\n'')

\section{7. Statistiques descriptives
------------------------------}\label{statistiques-descriptives}

stats \textless- df{[}, num\_vars{]} \%\textgreater\%
summarise(across(everything(), list( Min = \textasciitilde round(min(.),
3), Moy = \textasciitilde round(mean(.), 3), Med =
\textasciitilde round(median(.), 3), Max = \textasciitilde round(max(.),
3), SD = \textasciitilde round(sd(.), 3) ))) stats\_t \textless-
as.data.frame(t(stats))
stats\_t\(Variable <- gsub("_Min|_Moy|_Med|_Max|_SD", "", rownames(stats_t))
stats_t\)Stat \textless- gsub(``.*\_``,''``, rownames(stats\_t))
stats\_wide \textless- stats\_t \%\textgreater\%
pivot\_wider(names\_from = Stat, values\_from = V1) kable(stats\_wide,
caption =''Statistiques descriptives des variables numériques'',
booktabs = TRUE, digits = 3) \%\textgreater\%
kable\_styling(latex\_options = c(``striped'', ``hold\_position'',
``scale\_down''), font\_size = 9) cat(``\n\n'')

\section{8. Histogrammes (commentés)
-------------------------------}\label{histogrammes-commentuxe9s--}

\section{plots\_hist \textless- lapply(num\_vars, function(var)
\{}\label{plots_hist---lapplynum_vars-functionvar}

\section{ggplot(df, aes\_string(x = var))
+}\label{ggplotdf-aes_stringx-var}

\section{geom\_histogram(bins = 25, fill = ``\#2E86AB'', color =
``white'', alpha = 0.85)
+}\label{geom_histogrambins-25-fill-2e86ab-color-white-alpha-0.85}

\section{geom\_vline(aes(xintercept =
mean(df{[}{[}var{]}{]})),}\label{geom_vlineaesxintercept-meandfvar}

\section{color = ``black'', linetype = ``dashed'', linewidth = 0.8)
+}\label{color-black-linetype-dashed-linewidth-0.8}

\section{\texorpdfstring{labs(title = gsub(``\_``,'' ``, var), x =''``,
y =''Fréquence'')
+}{labs(title = gsub(``\_``,'' ``, var), x =''\,``, y =''Fréquence'') +}}\label{labstitle-gsub_-var-x-y-fruxe9quence}

\section{theme\_minimal(base\_size = 9)
+}\label{theme_minimalbase_size-9}

\section{theme(plot.title = element\_text(face = ``bold'', size =
9))}\label{themeplot.title-element_textface-bold-size-9}

\section{\})}\label{section-2}

\section{grid.arrange(grobs = plots\_hist, ncol =
3)}\label{grid.arrangegrobs-plots_hist-ncol-3}

\section{9. Matrice de corrélation (commentée)
---------------------}\label{matrice-de-corruxe9lation-commentuxe9e}

corr\_matrix \textless- cor(df{[}, num\_vars{]}) \#
corrplot(corr\_matrix, method = ``color'', type = ``lower'', \#
addCoef.col = ``black'', tl.col = ``black'', tl.srt = 45, \# number.cex
= 0.7, \# col = colorRampPalette(c(``\#C73E1D'', ``white'',
``\#2E86AB''))(200), \# mar = c(0, 0, 1, 0)) cat(``Corrélation
Happiness\_Score / GDP\_per\_Capita :'',
round(corr\_matrix{[}``Happiness\_Score'', ``GDP\_per\_Capita''{]}, 3),
``\n'') cat(``Corrélation Happiness\_Score / Healthy\_Life\_Expectancy
:'', round(corr\_matrix{[}``Happiness\_Score'',
``Healthy\_Life\_Expectancy''{]}, 3), ``\n'') cat(``Corrélation
Happiness\_Score / Education\_Index :'',
round(corr\_matrix{[}``Happiness\_Score'', ``Education\_Index''{]}, 3),
``\n'') cat(``Corrélation Happiness\_Score / Unemployment\_Rate :'',
round(corr\_matrix{[}``Happiness\_Score'', ``Unemployment\_Rate''{]},
3), ``\n\n'')

\section{10. Normalisation
-----------------------------------------}\label{normalisation}

X\_scaled \textless- scale(df{[}, num\_vars{]}) cat(``Vérification :
moyennes après normalisation (doivent être ≈ 0)\n'')
print(round(colMeans(X\_scaled), 5)) cat(``\n'')

\section{11. ACP
---------------------------------------------------}\label{acp}

res\_pca \textless- PCA(df{[}, num\_vars{]}, scale.unit = TRUE, graph =
FALSE)

\section{Variance expliquée}\label{variance-expliquuxe9e}

variance\_df \textless- data.frame( PC = paste0(``PC'', 1:9),
Valeur\_propre = round(res\_pca\(eig[, 1], 3),
  Variance_pct  = round(res_pca\)eig{[}, 2{]}, 2), Cumulative =
round(res\_pca\$eig{[}, 3{]}, 2) ) kable(variance\_df, caption =
``Variance expliquée par composante principale'', booktabs = TRUE,
col.names = c(``Composante'', ``Valeur propre (λ)'', ``Variance (\%)'',
``Cumulée (\%)'')) \%\textgreater\% kable\_styling(latex\_options =
c(``striped'', ``hold\_position'')) \%\textgreater\% row\_spec(1:2, bold
= TRUE, color = ``white'', background = ``\#1A5276'') cat(``\n'')

\section{Scree plot (commenté)}\label{scree-plot-commentuxe9}

\section{fviz\_eig(res\_pca, addlabels = TRUE, ylim = c(0,
75),}\label{fviz_eigres_pca-addlabels-true-ylim-c0-75}

\section{barfill = ``\#2E86AB'', barcolor = ``white'', linecolor =
``\#F18F01'') +}\label{barfill-2e86ab-barcolor-white-linecolor-f18f01}

\section{labs(title = ``Scree Plot --- Variance Expliquée par
Composante'')
+}\label{labstitle-scree-plot-variance-expliquuxe9e-par-composante}

\section{theme\_minimal(base\_size =
12)}\label{theme_minimalbase_size-12}

\section{Cercle des corrélations
(commenté)}\label{cercle-des-corruxe9lations-commentuxe9}

\section{fviz\_pca\_var(res\_pca, col.var =
``contrib'',}\label{fviz_pca_varres_pca-col.var-contrib}

\section{gradient.cols = c(``\#C73E1D'', ``\#F18F01'',
``\#2E86AB''),}\label{gradient.cols-cc73e1d-f18f01-2e86ab}

\section{repel = TRUE, title = ``Cercle des Corrélations (PC1 × PC2)'')
+}\label{repel-true-title-cercle-des-corruxe9lations-pc1-pc2}

\section{theme\_minimal(base\_size =
12)}\label{theme_minimalbase_size-12-1}

\section{Contributions (commentées)}\label{contributions-commentuxe9es}

\section{p\_c1 \textless- fviz\_contrib(res\_pca, choice = ``var'', axes
= 1, fill = ``\#2E86AB'', color = ``white'')
+}\label{p_c1---fviz_contribres_pca-choice-var-axes-1-fill-2e86ab-color-white}

\section{labs(title = ``Contributions à PC1'') +
theme\_minimal(base\_size =
10)}\label{labstitle-contributions-uxe0-pc1-theme_minimalbase_size-10}

\section{p\_c2 \textless- fviz\_contrib(res\_pca, choice = ``var'', axes
= 2, fill = ``\#F18F01'', color = ``white'')
+}\label{p_c2---fviz_contribres_pca-choice-var-axes-2-fill-f18f01-color-white}

\section{labs(title = ``Contributions à PC2'') +
theme\_minimal(base\_size =
10)}\label{labstitle-contributions-uxe0-pc2-theme_minimalbase_size-10}

\section{grid.arrange(p\_c1, p\_c2, ncol =
2)}\label{grid.arrangep_c1-p_c2-ncol-2}

\section{Ajout des coordonnées PC pour les couleurs dans le biplot
(commenté)}\label{ajout-des-coordonnuxe9es-pc-pour-les-couleurs-dans-le-biplot-commentuxe9}

df\(PC1 <- res_pca\)ind\(coord[, 1]
df\)PC2 \textless- res\_pca\(ind\)coord{[}, 2{]} \# couleurs\_region
\textless- c( \# ``Western Europe'' = ``\#2E86AB'', \# ``North America
\& ANZ'' = ``\#A23B72'', \# ``Latin America \& Caribbean'' =
``\#F18F01'', \# ``Central \& Eastern Europe'' = ``\#C73E1D'', \# ``East
Asia'' = ``\#44BBA4'', \# ``Southeast Asia'' = ``\#E94F37'', \#
``Sub-Saharan Africa'' = ``\#6C3483'' \# ) \#
fviz\_pca\_biplot(res\_pca, geom.ind = ``point'', col.ind = df\$Region,
\# palette = couleurs\_region, addEllipses = TRUE, \# ellipse.type =
``convex'', col.var = ``black'', repel = TRUE, \# legend.title =
``Région'', title = ``Biplot ACP : Individus et Variables'') + \#
theme\_minimal(base\_size = 11)

\section{12. Détermination du nombre de clusters (graphiques commentés)
----}\label{duxe9termination-du-nombre-de-clusters-graphiques-commentuxe9s--}

\section{fviz\_nbclust(X\_scaled, kmeans, method = ``wss'', k.max = 10,
linecolor = ``\#2E86AB'')
+}\label{fviz_nbclustx_scaled-kmeans-method-wss-k.max-10-linecolor-2e86ab}

\section{geom\_vline(xintercept = 4, linetype = ``dashed'', color =
``red'', linewidth = 1)
+}\label{geom_vlinexintercept-4-linetype-dashed-color-red-linewidth-1}

\section{labs(title = ``Méthode du Coude (Elbow)'', subtitle = ``Ligne
rouge = k=4 optimal'')
+}\label{labstitle-muxe9thode-du-coude-elbow-subtitle-ligne-rouge-k4-optimal}

\section{theme\_minimal(base\_size =
12)}\label{theme_minimalbase_size-12-2}

\section{fviz\_nbclust(X\_scaled, kmeans, method = ``silhouette'', k.max
= 10, linecolor = ``\#F18F01'')
+}\label{fviz_nbclustx_scaled-kmeans-method-silhouette-k.max-10-linecolor-f18f01}

\section{labs(title = ``Score de Silhouette par k'') +
theme\_minimal(base\_size =
12)}\label{labstitle-score-de-silhouette-par-k-theme_minimalbase_size-12}

\section{13. K-Means (k=4)
-----------------------------------------}\label{k-means-k4}

set.seed(42) km4 \textless- kmeans(X\_scaled, centers = 4, nstart = 25,
iter.max = 100) df\(Cluster <- as.factor(km4\)cluster) cat(``=== Taille
des clusters ===\n'') print(table(df\(Cluster))
sil <- silhouette(km4\)cluster, dist(X\_scaled))
cat(sprintf(``\nScore de silhouette moyen : \%.4f\n'', mean(sil{[},
3{]}))) cat(sprintf(``Inertie expliquée : \%.2f\%\%\n'', 100 *
km4\(betweenss / km4\)totss)) cat(``\n'')

\section{Visualisation des clusters
(commentée)}\label{visualisation-des-clusters-commentuxe9e}

\section{fviz\_cluster(km4, data = X\_scaled, ellipse.type =
``convex'',}\label{fviz_clusterkm4-data-x_scaled-ellipse.type-convex}

\section{palette = c(``\#2E86AB'', ``\#F18F01'', ``\#C73E1D'',
``\#44BBA4''),}\label{palette-c2e86ab-f18f01-c73e1d-44bba4}

\section{ggtheme = theme\_minimal(base\_size =
12),}\label{ggtheme-theme_minimalbase_size-12}

\section{main = ``Clusters K-Means (k=4) dans l'Espace ACP'', repel =
TRUE)}\label{main-clusters-k-means-k4-dans-lespace-acp-repel-true}

\section{fviz\_silhouette(sil, palette = c(``\#2E86AB'', ``\#F18F01'',
``\#C73E1D'',
``\#44BBA4''),}\label{fviz_silhouettesil-palette-c2e86ab-f18f01-c73e1d-44bba4}

\section{ggtheme = theme\_minimal(base\_size =
11),}\label{ggtheme-theme_minimalbase_size-11}

\section{main = ``Diagramme de Silhouette
(k=4)'')}\label{main-diagramme-de-silhouette-k4}

\section{14. CAH (dendrogramme commenté)
---------------------------}\label{cah-dendrogramme-commentuxe9}

set.seed(42) sample\_idx \textless- sample(1:nrow(X\_scaled), 80)
dist\_mat \textless- dist(X\_scaled{[}sample\_idx, {]}, method =
``euclidean'') hc\_ward \textless- hclust(dist\_mat, method =
``ward.D2'') dend \textless- as.dendrogram(hc\_ward) dend \textless-
color\_branches(dend, k = 4, col = c(``\#2E86AB'', ``\#F18F01'',
``\#C73E1D'', ``\#44BBA4'')) dend \textless- set(dend, ``labels\_cex'',
0.45) \# plot(dend, main = ``Dendrogramme CAH (Ward, n=80)'', ylab =
``Distance de Ward'', xlab = ``Observations'') \# abline(h = 8, col =
``red'', lty = 2, lwd = 2) \# legend(``topright'', legend = ``Seuil de
coupure (4 clusters)'', lty = 2, col = ``red'', lwd = 2, cex = 0.85)

\section{15. Profils moyens des clusters
---------------------------}\label{profils-moyens-des-clusters}

profils \textless- df \%\textgreater\% group\_by(Cluster)
\%\textgreater\% summarise(across(all\_of(num\_vars),
\textasciitilde round(mean(.), 3)), n = n()) kable(t(profils), caption =
``Profils moyens des 4 clusters K-Means'', booktabs = TRUE)
\%\textgreater\% kable\_styling(latex\_options = c(``striped'',
``hold\_position'', ``scale\_down''), font\_size = 9) cat(``\n'')

\section{Heatmap des profils
(commenté)}\label{heatmap-des-profils-commentuxe9}

\section{profils\_long \textless- profils
\%\textgreater\%}\label{profils_long---profils}

\section{select(-n) \%\textgreater\%}\label{select-n}

\section{pivot\_longer(-Cluster, names\_to = ``Variable'', values\_to =
``Valeur'')
\%\textgreater\%}\label{pivot_longer-cluster-names_to-variable-values_to-valeur}

\section{group\_by(Variable) \%\textgreater\%}\label{group_byvariable}

\section{mutate(Valeur\_norm = (Valeur - min(Valeur)) / (max(Valeur) -
min(Valeur)))}\label{mutatevaleur_norm-valeur---minvaleur-maxvaleur---minvaleur}

\section{ggplot(profils\_long, aes(x = Cluster, y = Variable, fill =
Valeur\_norm))
+}\label{ggplotprofils_long-aesx-cluster-y-variable-fill-valeur_norm}

\section{geom\_tile(color = ``white'', linewidth = 0.5)
+}\label{geom_tilecolor-white-linewidth-0.5}

\section{geom\_text(aes(label = round(Valeur, 2)), size = 3, color =
``black'') +}\label{geom_textaeslabel-roundvaleur-2-size-3-color-black}

\section{scale\_fill\_gradient2(low = ``\#C73E1D'', mid = ``white'',
high =
``\#2E86AB'',}\label{scale_fill_gradient2low-c73e1d-mid-white-high-2e86ab}

\section{\texorpdfstring{midpoint = 0.5, name = ``Valeur\nnormalisée'')
+}{midpoint = 0.5, name = ``Valeur'') +}}\label{midpoint-0.5-name-valeur}

\section{labs(title = ``Heatmap des Profils de Clusters'', x =
``Cluster'', y = ``\,``)
+}\label{labstitle-heatmap-des-profils-de-clusters-x-cluster-y}

\section{theme\_minimal(base\_size =
12)}\label{theme_minimalbase_size-12-3}

\section{Composition régionale des clusters (graphique
commenté)}\label{composition-ruxe9gionale-des-clusters-graphique-commentuxe9}

\section{cross\_tab \textless- df
\%\textgreater\%}\label{cross_tab---df}

\section{group\_by(Cluster, Region)
\%\textgreater\%}\label{group_bycluster-region}

\section{summarise(n = n(), .groups = ``drop'')
\%\textgreater\%}\label{summarisen-n-.groups-drop}

\section{group\_by(Cluster) \%\textgreater\%}\label{group_bycluster}

\section{mutate(pct = round(100 * n / sum(n),
1))}\label{mutatepct-round100-n-sumn-1}

\section{ggplot(cross\_tab, aes(x = Cluster, y = pct, fill = Region))
+}\label{ggplotcross_tab-aesx-cluster-y-pct-fill-region}

\section{geom\_bar(stat = ``identity'', color = ``white'')
+}\label{geom_barstat-identity-color-white}

\section{scale\_fill\_manual(values = couleurs\_region)
+}\label{scale_fill_manualvalues-couleurs_region}

\section{labs(title = ``Composition Régionale des Clusters'', x =
``Cluster'', y = ``Proportion (\%)'', fill = ``Région'')
+}\label{labstitle-composition-ruxe9gionale-des-clusters-x-cluster-y-proportion-fill-ruxe9gion}

\section{theme\_minimal(base\_size =
12)}\label{theme_minimalbase_size-12-4}

\section{16. Interprétation
----------------------------------------}\label{interpruxe9tation--}

best \textless- profils\(Cluster[which.max(profils\)Happiness\_Score){]}
cat(sprintf(``Cluster le plus heureux : \%s \textbar{} Bonheur=\%.2f
\textbar{} PIB=\%.2f \textbar{} Education=\%.2f \textbar{}
Chômage=\%.1f\%\%\n'', best, profils\(Happiness_Score[profils\)Cluster
== best{]}, profils\(GDP_per_Capita[profils\)Cluster == best{]},
profils\(Education_Index[profils\)Cluster == best{]},
profils\(Unemployment_Rate[profils\)Cluster == best{]}))

worst \textless-
profils\(Cluster[which.min(profils\)Happiness\_Score){]}
cat(sprintf(``Cluster le moins heureux : \%s \textbar{} Bonheur=\%.2f
\textbar{} PIB=\%.2f \textbar{} Chômage=\%.1f\%\%\n'', worst,
profils\(Happiness_Score[profils\)Cluster == worst{]},
profils\(GDP_per_Capita[profils\)Cluster == worst{]},
profils\(Unemployment_Rate[profils\)Cluster == worst{]}))

\section{17. Résumé final
------------------------------------------}\label{ruxe9sumuxe9-final}

cat(``\n=============================================================\n'')
cat('' RÉSUMÉ DE L'ANALYSE MULTIVARIÉE\n``)
cat(''=============================================================\n``)
cat(sprintf('' Dataset : \%d pays x \%d variables\n``, nrow(df),
length(num\_vars))) cat(sprintf('' PC1 variance : \%.2f\%\%\n``,
res\_pca\(eig[1, 2]))
cat(sprintf("  PC1+PC2 variance : %.2f%%\n", res_pca
\)eig{[}2, 3{]})) cat(sprintf('' Clusters K-Means : k = 4\n``))
cat(sprintf('' Score silhouette : \%.4f\n``, mean(sil{[}, 3{]})))
cat(sprintf('' Inertie expliquée: \%.2f\%\%\n``, 100 *
km4\(betweenss / km4\)totss))
cat(''=============================================================\n``)

\end{document}
